\clearpage\SourceLeaf{299}
\begin{center}
\textbf{[GEOMETRISCHER ORT}\par\medskip
\textbf{DER SPITZE DES SPHÄRISCHEN DREIECKS}\par\medskip
\textbf{AUF GEGEBENER BASIS MIT GEGEBENEM INHALT.]}
\end{center}

\medskip

\begin{center}
\textsc{Gauss} an \textsc{Schumacher}. Göttingen, 29. December 1841.
\end{center}

\mbox{.\,.\,.\,.\,.\,.\,.} Ist Ihnen eine \textls{gedruckte} Auflösung der Aufgabe: den geometrischen Ort der Spitze eines sphärischen Dreiecks auf gegebener Basis und von gegebenem Inhalt zu finden, bekannt? Ich wurde neulich veranlasst sie zu suchen; und war etwas verwundert, dass die analytische Auflösung mit einigem Umwege auf ein höchst einfaches Resultat führt, dessen synthetischer Beweis sich in ein Paar Zeilen führen lässt. \mbox{.\,.\,.\,.\,.\,.\,.}

\medskip

\begin{center}
[\textsc{Schumacher} \emph{an} \textsc{Gauss}.\ \emph{Altona, 3. Januar 1842.}]
\end{center}

\{Eine geometrische Auflösung des von Ihnen, mein theuerster Freund, erwähnten Problems hat vor nicht gar langer Zeit \textsc{Steiner} in \textsc{Crelles} Journal (Bd.\ 2.\ p.\ 45) gegeben, wo er aber seine Vorgänger nicht richtig citirt. \textsc{Lexells} Aufsatz steht nicht, wie dort gesagt wird, in Nova Acta Petrop.\ V.\ Pars 1, sondern in Acta Petrop.\ pro 1781, Pars I. \textsc{Euler} ist später darauf in Nova Acta Tom.\ X.\ p.\ 47 (Variae speculationes etc.) zurückgekommen. \textsc{Legendre} hat ihn
